\subsection{The Cosine Theorem Emerges Immediately}

Draw $\mathbf u$ and $\mathbf v$ from a common point. The displacement from the
endpoint of $\mathbf v$ to the endpoint of $\mathbf u$ is
\[
\mathbf u-\mathbf v.
\]
Thus these three vectors determine a triangle with side lengths
\[
\|\mathbf u\|,
\qquad
\|\mathbf v\|,
\qquad
\|\mathbf u-\mathbf v\|.
\]

\begin{derivationbox}[Expanding the third side]
\[
\begin{aligned}
\|\mathbf u-\mathbf v\|^2
&=(\mathbf u-\mathbf v)\cdot(\mathbf u-\mathbf v)\\
&=\mathbf u\cdot\mathbf u
-\mathbf u\cdot\mathbf v
-\mathbf v\cdot\mathbf u
+\mathbf v\cdot\mathbf v\\
&=\|\mathbf u\|^2+\|\mathbf v\|^2-2\mathbf u\cdot\mathbf v\\
&=\|\mathbf u\|^2+\|\mathbf v\|^2
-2\|\mathbf u\|\,\|\mathbf v\|\cos\theta.
\end{aligned}
\]
If
\[
a=\|\mathbf u\|,
\qquad
b=\|\mathbf v\|,
\qquad
c=\|\mathbf u-\mathbf v\|,
\]
then the familiar triangle formula follows.
\end{derivationbox}

\begin{theorembox}
For a triangle with side lengths $a$, $b$, and $c$, where $\theta$ is the angle
between the sides of lengths $a$ and $b$,
\[
\boxed{c^2=a^2+b^2-2ab\cos\theta}.
\]
\end{theorembox}

\begin{interpretationbox}[Why the theorem appears here]
The cosine theorem is not inserted as an external formula. It is the expansion
of the squared displacement between the endpoints of two arbitrary vectors.
The triangle law is therefore encoded in vector subtraction and the dot product.
\end{interpretationbox}

\begin{center}
\begin{tikzpicture}[scale=0.9]
    \coordinate (O) at (0,0);
    \coordinate (U) at (4.4,1.5);
    \coordinate (V) at (1.3,3.5);
    \draw[subset,->] (O) -- (U) node[midway,below] {$\mathbf u$};
    \draw[subset,->] (O) -- (V) node[midway,left] {$\mathbf v$};
    \draw[very thick,->] (V) -- (U) node[midway,above right] {$\mathbf u-\mathbf v$};
    \draw (0.72,0.25) arc[start angle=19,end angle=69.6,radius=0.76];
    \node at (0.70,0.65) {$\theta$};
\end{tikzpicture}
\end{center}
